\chapter{Experimentos}
\label{cap:experimentos}

Este capítulo distingue dos protocolos con preguntas y unidades experimentales diferentes. El \emph{benchmark legado} sobre cinco series reales, ya ejecutado, usa historias y horizontes medidos en eventos. El \emph{benchmark físico sintético} usa duraciones continuas, targets extraídos de la trayectoria latente limpia con independencia del patrón de observación y del ruido, y controles de identificabilidad; su ranking neuronal sólo se incorpora cuando termina la ejecución completa. Los scores de ambos protocolos no son directamente comparables.

\section{Objetivos experimentales}
\label{sec:objetivos-experimentales}

Los experimentos persiguen los siguientes objetivos:

\begin{enumerate}
	\item \textbf{Verificar factibilidad}: comprobar que arquitecturas basadas en \Transformers\ con timestamps reales y tokens objetivo pueden entrenarse de manera estable sobre series reales no equiespaciadas sin regularizarlas a una grilla temporal fija.
	\item \textbf{Evaluar desempeño predictivo}: medir MSE, RMSE y MAE en el conjunto de test retenido.
	\item \textbf{Comparar contra baselines informados}: contrastar la propuesta con STraTS \cite{tipirneni2022strats} y CoFormer \cite{wei2023coformer}, adaptados a la tarea de pronóstico multi-horizonte.
	\item \textbf{Comparar contra baselines simples}: incluir persistencia y un modelo lineal temporal para verificar que las mejoras no provengan sólo de aumentar la complejidad.
	\item \textbf{Cuantificar ablaciones}: medir el aporte del encoding temporal continuo y del identificador de token objetivo.
	\item \textbf{Explorar variantes arquitectónicas}: evaluar tamaños pequeños y base, resumen atencional de la historia, Time2Vec, sesgo temporal de atención y variantes encoder--decoder.
	\item \textbf{Evaluar costo versus beneficio}: registrar parámetros, tiempo de entrenamiento y épocas ejecutadas.
	\item \textbf{Delimitar el alcance de las conclusiones}: separar las señales consistentes del benchmark cerrado de las afirmaciones que requerirían más datasets, más semillas o evaluación multivariada.
\end{enumerate}

\section{Benchmark legado sobre series reales}
\label{sec:escenario-real}

La evaluación legada se realiza sobre los archivos \textttbreak{real\_data\_i.csv}, con $i=1,\dots,5$. Cada dataset se ordena temporalmente y se divide respetando el tiempo en entrenamiento, validación y prueba, con proporciones 70\%, 15\% y 15\%, respectivamente. Para cada modelo se ejecutan tres semillas: 42, 84 y 126.

\hfill

El archivo principal de resultados es \textttbreak{benchmark\_final.csv}, ubicado en el directorio de experimentos finales. El subconjunto usado para el análisis final contiene 360 corridas balanceadas: 24 modelos con 15 corridas cada uno, correspondientes a 5 datasets y 3 semillas. Las filas parciales asociadas a otros datasets o variantes incompletas no se usan para extraer conclusiones finales.

\hfill

El benchmark se ejecuta mediante un orquestador reanudable. Después de cada corrida, las métricas se guardan en un archivo CSV; si el proceso se interrumpe, puede continuar desde el último punto completado. Esta decisión fue importante porque algunos modelos, en especial CoFormer y las variantes encoder--decoder con ajuste autoregresivo, tienen costos de entrenamiento sustanciales.

\subsection{Variables, partición temporal y preprocesamiento}
\label{sec:variables-y-preproceso}

En la configuración activa de datos reales, cada serie contiene:
\begin{itemize}
	\item una columna temporal \textttbreak{timestamp},
	\item una variable observada \textttbreak{valor}, usada como entrada y como objetivo de predicción,
	\item una partición temporal entrenamiento--validación--prueba de 70\%--15\%--15\%.
\end{itemize}

El escalado se ajusta sólo con el segmento de entrenamiento y luego se aplica a validación y prueba, evitando fuga de información desde el futuro. Las ventanas supervisadas usan una historia máxima de 500 observaciones y una historia mínima de 20 observaciones. Para los modelos no autoregresivos se seleccionan hasta cuatro timestamps objetivo dentro de offsets entre 1 y 30 pasos futuros. La escala temporal del modelo se estima por dataset a partir de la mediana de los deltas temporales positivos del bloque de entrenamiento, con límites conservadores para evitar escalas patológicas.

\section{Modelos comparados en el protocolo legado}
\label{sec:modelos-comparados}

Los modelos evaluados se agrupan en cinco familias.

\subsection{Familia \CTTransformer}

\begin{itemize}
	\item \textbf{\CTTransformerBase}: encoder \Transformer\ con embedding de valores, encoding temporal sinusoidal continuo, identificador historia--objetivo y cabeza de regresión sobre los tokens objetivo.
	\item \textbf{\CTTransformerSmall}: misma arquitectura, con menor dimensión interna y menos capas.
	\item \textbf{\CTTransformer\ con resumen atencional}: fusiona el estado del token objetivo con un resumen aprendido de la historia.
	\item \textbf{\CTTransformer\ con Time2Vec}: reemplaza el encoding sinusoidal fijo por una representación temporal aprendible basada en Time2Vec \cite{kazemi2019time2vec}.
	\item \textbf{\CTTransformer\ con sesgo temporal}: conserva el encoding sinusoidal, pero añade un sesgo temporal directamente en los scores de atención.
	\item \textbf{\CTTransformer\ con Time2Vec y sesgo temporal}: combina representación temporal aprendida y sesgo temporal de atención.
\end{itemize}

\subsection{Variantes encoder--decoder}

Además del encoder único, se evalúan variantes encoder--decoder donde la historia se procesa en el encoder y los timestamps objetivo se procesan en el decoder mediante atención cruzada. La versión optimizada usa un decoder liviano de una capa. Se consideran dos tamaños: pequeño y base.

\hfill

Las variantes \textttbreak{EncDec-Opt-Small-MT8} y \textttbreak{EncDec-Opt-MT8} se entrenan con ocho tokens objetivo. Después se aplican ajustes autoregresivos con máscara causal y tres modos de provisión de objetivos:
\begin{itemize}
	\item \textbf{Contiguous}: offsets futuros consecutivos.
	\item \textbf{Random}: offsets futuros muestreados desde el rango permitido.
	\item \textbf{Mixed}: mezcla entre entrenamiento contiguo y aleatorio.
\end{itemize}

En las variantes autoregresivas, la evaluación registrada usa generación autoregresiva sobre ocho objetivos contiguos. Por esta razón, cuando se compara contra modelos no autoregresivos, la lectura correcta es que se compara el desempeño de los sistemas tal como fueron evaluados en el benchmark, no una demostración aislada bajo exactamente el mismo conjunto de consultas futuras.

\subsection{Baselines del estado del arte}

STraTS y CoFormer se incluyen porque son dos referencias representativas para series irregulares sin imputación a una malla fija. No obstante, ambos fueron propuestos principalmente como codificadores de observaciones irregulares para tareas de clasificación, por lo que en este benchmark se usan adaptaciones orientadas a pronóstico multi-horizonte.

\begin{itemize}
	\item \textbf{STraTS adaptado}: conserva la representación por tripletas tiempo--variable--valor y el encoder con fusión por atención de STraTS. Para hacerlo comparable con la tarea de esta memoria, se añaden consultas futuras marcadas como objetivos, de modo que el modelo pueda entregar una predicción por timestamp objetivo. Esta adaptación no convierte a STraTS en un encoder--decoder: sigue siendo un codificador de tripletas con lectura sobre representaciones contextualizadas.
	\item \textbf{CoFormer adaptado}: conserva la idea de puntos variable--tiempo y la separación entre atención intra-variable e inter-variable. En la configuración actual de datos reales, las corridas consolidadas son univariadas, por ello CoFormer opera como \textttbreak{CoFormer-Uni}, donde la atención inter-variable deja de aportar el mecanismo central para el que fue diseñado. Para pronóstico, la representación agregada de la historia se combina con la información temporal de cada objetivo futuro.
\end{itemize}

\subsection{Baselines simples}

\begin{itemize}
	\item \textbf{Persistencia}: predice el último valor observado de la historia para todos los horizontes futuros.
	\item \textbf{Modelo lineal temporal}: usa los últimos 50 valores de la historia junto con tiempos relativos para predecir cada horizonte.
\end{itemize}

\subsection{Ablaciones}

\begin{itemize}
	\item \textbf{Sin encoding temporal continuo}: reemplaza los timestamps reales por un encoding posicional ordinal. Mide el aporte específico del tiempo continuo.
	\item \textbf{Sin identificador de token objetivo}: elimina la señal aprendible que distingue historia y objetivo. Mide el aporte de explicitar el rol del token.
\end{itemize}

\section{Configuración de entrenamiento legada}
\label{sec:configuracion-entrenamiento}

Las configuraciones paramétricas principales son:
\begin{itemize}
	\item \textbf{Pequeña}: $d_{\text{model}} = 64$, 2 capas, 4 cabezas, \emph{feed-forward} 128, \emph{dropout} 0.1.
	\item \textbf{Base}: $d_{\text{model}} = 128$, 4 capas, 4 cabezas, \emph{feed-forward} 512, \emph{dropout} 0.15.
\end{itemize}

El entrenamiento usa AdamW, calentamiento lineal seguido de decaimiento coseno, selección del mejor checkpoint según RMSE de validación, precisión mixta automática y recorte de gradientes. La configuración se ajusta según el tamaño del modelo. En las variantes más inestables de \CTTransformer\ y del encoder--decoder optimizado se usa pérdida Huber, menor tasa de aprendizaje, más calentamiento y recorte de gradiente más estricto. El scheduler se detalla en el Anexo~\ref{anexo:optimizacion-scheduler}.

\section{Métricas de evaluación legadas}
\label{sec:metricas-evaluacion-experimentos}

La evaluación se realiza mediante:
\begin{itemize}
	\item \textbf{MSE}: métrica principal, sensible a errores grandes.
	\item \textbf{RMSE}: raíz del MSE, expresada en la misma escala de la variable normalizada.
	\item \textbf{MAE}: error absoluto medio, menos sensible a valores atípicos.
	\item \textbf{MAPE}: error porcentual absoluto medio, se reporta con cautela porque puede ser inestable cuando el valor real está cerca de cero.
\end{itemize}

El análisis principal prioriza MSE, RMSE y MAE.

\section{Agregación y análisis estadístico del benchmark legado}
\label{sec:analisis-estadistico-experimentos}

Para cada modelo, las tres semillas se promedian primero dentro de cada dataset. Luego se reporta la media y desviación estándar sobre los cinco promedios por dataset. Esta decisión evita tratar semillas como si fueran datasets independientes y permite que cada serie tenga el mismo peso en las tablas finales.

\hfill

También se reportan victorias por dataset, calculadas con MSE luego de promediar semillas. Dado que sólo hay cinco datasets, las pruebas estadísticas emparejadas tienen baja potencia. Por ello, las conclusiones del Capítulo~\ref{cap:resultados} se basan en magnitud, consistencia entre datasets y comparación de costos, no en una afirmación de significancia estadística fuerte.

\section{Costo computacional}
\label{sec:costo-computacional-experimentos}

Para cada corrida se registra el número de parámetros entrenables, el número total de parámetros, el tiempo de entrenamiento en segundos y el número de épocas ejecutadas antes de detener el entrenamiento. En las variantes con ajuste autoregresivo, el costo reportado en las tablas finales corresponde al entrenamiento base más el ajuste posterior, porque ambos pasos son necesarios para obtener el modelo evaluado.

\section{Reproducibilidad}
\label{sec:reproducibilidad}

El benchmark se puede reanudar ejecutando:
\begin{verbatim}
python scripts/benchmark_final.py
\end{verbatim}

El orquestador admite seleccionar rangos de datasets, semillas, subconjuntos de modelos y variantes opcionales como encoder--decoder y ajuste autoregresivo. Para reproducir el análisis final de esta memoria se debe filtrar el archivo de resultados a los datasets 1--5, semillas 42, 84 y 126, y modelos con las 15 corridas completas.

\section{Evaluación física sintética congelada}
\label{sec:benchmark-fisico}

La segunda evaluación está especificada en \nolinkurl{configs/benchmark/thesis_physical_final.yaml} y se ejecuta exclusivamente mediante \nolinkurl{scripts/run_thesis_physical_benchmark.py}. Su estado versionado exacto es \nolinkurl{thesis_physical_evaluation_v3_frozen}. El archivo YAML fija a priori datos, semillas, modelos, presupuesto, selección de checkpoints, métricas, política de ejecución y rutas de salida. El runner no ejecuta tuning automático ni modifica hiperparámetros a partir del conjunto de test. Su interpretación es retrospectiva: los procesos sintéticos ya generados y sus diagnósticos fueron inspeccionados durante el desarrollo, por lo que no se presenta como una prueba prospectiva ciega.

\subsection{Unidades experimentales y partición principal--estrés}

La partición principal contiene exactamente 16 unidades: los identificadores canónicos \textttbreak{*\_0000} de ocho presets sintéticos, cada uno en modalidad univariada y multivariada. Se excluyen del estimando principal las dos unidades \textttbreak{long\_gaps\_gseed3031\_0000}, una por modalidad, que forman una partición de estrés separada. Esta separación evita incorporar a posteriori un caso extremo al promedio principal y permite reportarlo como prueba de sensibilidad.

Cada unidad provee \nolinkurl{observations.parquet}, con eventos irregulares y ruido, y \nolinkurl{truth.parquet}, con la trayectoria limpia. Los splits temporales se respetan de extremo a extremo. El escalado se ajusta sólo en entrenamiento y los targets se obtienen por interpolación lineal de \nolinkurl{truth.parquet}; test no participa en tuning, parada temprana ni selección del checkpoint.

\subsection{Tarea en tiempo físico}

Todas las historias cubren una duración física de 8.0 unidades. Para limitar memoria se conservan como máximo 256 eventos en univariado y 512 en multivariado mediante submuestreo representativo en tiempo. Los orígenes de forecast se toman de una grilla física de la trayectoria limpia, no de la llegada de un evento. Una historia sin observaciones se descarta y cada descarte se asigna a una causa mutuamente excluyente.

Durante entrenamiento cada ejemplo usa cuatro horizontes muestreados con una distribución log-uniforme en $[0.25,8.0]$. La evaluación usa $\{0.25,1.0,3.0,8.0\}$ y permuta reproduciblemente el orden de las consultas. Los timestamps absolutos permanecen en \textttbreak{float64} durante carga, selección y resta del origen; las capas reciben sólo tiempos relativos ya recentrados. Una auditoría explícita compara esta ruta con el cast absoluto prematuro a \textttbreak{float32}.

\subsection{Modelos, ablaciones y selección}

Se ejecutan las semillas 42, 84 y 126 para los siguientes sistemas:
\begin{itemize}
    \item \textttbreak{QueryCross} y \textttbreak{BasisDecoder} como arquitecturas puntuales principales;
    \item \textttbreak{CTSSM} y \textttbreak{BasisDecoder-CTSSM}, que añaden la transición continua estable;
    \item \textttbreak{NoTime} y \textttbreak{QueryOnly} como ablaciones temporales de \QueryCross;
    \item \textttbreak{QueryCross-Gaussian} y \textttbreak{BasisDecoder-Gaussian} como extensiones probabilísticas;
    \item \textttbreak{Persistence} como baseline sin parámetros.
\end{itemize}
Las variantes puntuales minimizan su pérdida de regresión y seleccionan checkpoint por \textttbreak{val\_rmse}. Las gaussianas minimizan NLL heteroscedástica y seleccionan por \textttbreak{val\_nll}. La receta física usa AdamW y \textttbreak{CosineAnnealingLR} sin warmup, a diferencia del warmup más coseno del benchmark legado. Las comparaciones de ablación reutilizan inicializaciones emparejadas cuando el grafo común lo permite. Los hiperparámetros permanecen fijados en el YAML antes de evaluar test. La campaña usa batch 32, \textttbreak{num\_workers=0} y CUDA obligatoria; el preflight verifica disponibilidad de la GPU RTX 5090 y memoria antes de habilitar la corrida completa.

La versión v3 activa TF32 y los kernels SDPA rápidos disponibles, incluidos Flash Attention o la ruta eficiente en memoria cuando PyTorch puede utilizarlos. Las seeds siguen fijando muestreo, inicialización y orden de ejemplos, pero estos kernels no garantizan igualdad bit a bit entre ejecuciones o versiones de CUDA. En consecuencia, la reproducibilidad reclamada es estadística y trazable: se conservan tres seeds, configuración efectiva, hashes de código y datos, entorno numérico y métricas completas; no se afirma determinismo binario.

\subsection{Identificabilidad y métricas}

Antes del ranking neuronal se audita que el slot de query no determine el horizonte y se comparan controles de persistencia, representación ordinal y horizonte explícito. Las corruptelas emparejadas incluyen timestamps reales, igualados, gaps permutados, grilla regular, orden ordinal, \textttbreak{query\_only} y \textttbreak{history\_only}. Estos diagnósticos distinguen sensibilidad temporal de una simple asociación con orden o densidad; no demuestran por sí solos causalidad sobre datos reales.

Para predicción puntual se reportan RMSE, MAE y sesgo en espacio estandarizado y original. Para cabezas gaussianas se añaden NLL, CRPS, escala predictiva y cobertura de intervalos de 90\% y 95\%. Las métricas se estratifican por horizonte, densidad, hueco máximo y edad de la última observación; en multivariado los cuantiles se calculan dentro de cada canal para no confundir identidad del sensor con irregularidad.

La agregación respeta la jerarquía experimental. Primero se promedian las semillas dentro de cada unidad dataset; luego se emparejan las modalidades univariada y multivariada dentro de cada preset; finalmente se calcula el macro-promedio entre los ocho presets. Las semillas miden variabilidad de optimización y no aumentan el número de réplicas científicas. El stress set se informa aparte y nunca entra al ranking principal.

\subsection{Ejecución y estado de los resultados}

El runner ofrece los modos \textttbreak{preflight}, \textttbreak{dry-run}, \textttbreak{smoke}, \textttbreak{run} y \textttbreak{report}, además de \textttbreak{all} para ejecutar la secuencia completa. El modo \textttbreak{smoke} usa un presupuesto reducido para validar lectura de datos, forward, backward, checkpoints y reportes; sus métricas no constituyen un ranking. Sólo \textttbreak{run} completo, seguido de \textttbreak{report}, puede poblar las tablas finales del Capítulo~\ref{cap:resultados}. La salida v3 se guarda bajo \nolinkurl{experiments/thesis_physical_benchmark_v3/}; los artefactos de v1 y v2 permanecen aislados y no se reutilizan.

Cada cohorte usa un único \textttbreak{shard\_0}. El runner completa primero la cohorte principal y luego la de estrés, con un solo proceso de entrenamiento residente en la GPU; dentro de él, cada combinación de dataset, modelo y seed se ejecuta secuencialmente. Cada dataset conserva un \textttbreak{protocol\_index} canónico, por lo que esta política no modifica la seed usada para construir historias y queries. El proceso escribe logs y sentinels propios, y sólo se publica un consolidado después de validar la cobertura exacta sin duplicados ni faltantes. La lectura de Parquet se omite al reanudar un dataset ya completo y compatible. Durante evaluación se reutilizan transferencias CPU--GPU, se vectoriza la construcción de métricas y se preindexan diagnósticos por sensor. La historia determinista submuestreada puede cachearse; los targets nunca se cachean y las queries de entrenamiento continúan cambiando por época.

La política secuencial se eligió mediante probes de ingeniería ajenos al ranking científico. Sobre el mismo conjunto de cuatro trabajos, cuatro procesos concurrentes lograron mayor throughput que dos: 35.928 segundos frente a 57.278. Esa alternativa no se adoptó en v3 porque aumenta el uso y la contención de CPU, mantiene varios procesos y modelos activos a la vez y complica la observación y el reinicio de una campaña ejecutada sobre una sola GPU. La política seleccionada prioriza un consumo predecible y una operación sencilla; no se afirma que minimice el tiempo de pared absoluto. También se descartaron batches mayores y workers adicionales: el primero alteraba la dinámica de optimización y el segundo no redujo el tiempo total en Windows. Estas mediciones justifican la configuración operativa, pero no se interpretan como resultados de precisión del modelo.

La ejecución completa se inicia con un único comando:
\begin{verbatim}
conda.exe run -n memoria python `
  scripts/run_thesis_physical_benchmark.py all `
  --config configs/benchmark/thesis_physical_final.yaml
\end{verbatim}

El subcomando \textttbreak{all} ejecuta preflight, las cohortes físicas en secuencia, controles de identificabilidad y reporte. Ante una interrupción, repetir el mismo comando reanuda únicamente artefactos completos con fingerprint compatible.

\section{Resumen metodológico}
\label{sec:resumen-metodologico}

En síntesis, el benchmark legado compara \CTTransformer, variantes encoder--decoder, baselines simples y adaptaciones de STraTS y CoFormer sobre cinco series reales; sus resultados son descriptivos a nivel de dataset. La evaluación física congelada compara \QueryCross, \ContinuousBasis, \CTSSM, ablaciones y cabezas gaussianas sobre 16 unidades sintéticas principales y dos de estrés, con tiempo identificable y agregación jerárquica. El capítulo siguiente conserva ambos cuerpos de evidencia separados y no completa las tablas físicas con resultados de smoke tests.
